\section{Materiales y métodos}

\subsection{Estado actual y brecha de conocimiento}

Diversos estudios han explorado los factores genéticos y fisiológicos implicados en los comportamientos anómalos de ingesta, especialmente en relación con el consumo de alcohol. Morozova et al. (2012) \cite{morozova2012} describieron la complejidad genética del alcoholismo, identificando múltiples genes y fenotipos asociados, mientras que Rachdaoui y Sarkar (2013) \cite{rachdaoui2013} detallaron cómo el alcohol altera los mecanismos endocrinos que regulan la homeostasis hídrica. Sin embargo, ambos enfoques se centran en la respuesta al alcohol y no en las alteraciones primarias del impulso de beber o en la regulación molecular de la sed.

En el ámbito clínico, Bhatia y Sharma (2017) \cite{bhatia2017} abordaron la polidipsia psicogénica y sus dificultades de manejo, y Yamashiro et al. (2013) \cite{yamashiro2013} documentaron casos de intoxicación por agua. Aunque estas investigaciones evidencian la relevancia fisiológica y médica del comportamiento de ingesta anómala, se limitan a descripciones clínicas y carecen de una integración con modelos moleculares o genómicos.

Más recientemente, Leger et al. (2024) \cite{leger2024} aplicaron un enfoque de biología de sistemas para analizar redes génicas asociadas al consumo de alcohol, demostrando el potencial de estos métodos para comprender comportamientos complejos. No obstante, no existen estudios equivalentes que apliquen este enfoque al fenotipo \textit{Abnormal drinking behaviour} (HP:0030082), tal como se define en la \textbf{Human Phenotype Ontology} (HPO). Por tanto, la principal brecha en el conocimiento reside en la falta de un análisis sistemático de redes génicas que vincule este fenotipo con los mecanismos moleculares que regulan la homeostasis hídrica y el control de la sed.

\subsection{Hipótesis de trabajo}

Se plantea la hipótesis de que los genes asociados al fenotipo \textit{Abnormal drinking behaviour} (HP:0030082) presentan una interconexión funcional y participan en rutas moleculares vinculadas con la regulación de la homeostasis hídrica y el control neuroendocrino de la sed. 

En consecuencia, el análisis de sus interacciones mediante enfoques de \textbf{biología de sistemas} —por ejemplo, a través de la plataforma \textit{STRING}— permitirá identificar módulos funcionales y genes clave (\textit{hubs}) que podrían constituir la base molecular común de los comportamientos anómalos de ingesta descritos en la literatura clínica.

\subsection{Objetivo general}

Investigar cómo se relacionan los genes asociados al fenotipo \textit{Abnormal drinking behaviour} (HP:0030082) utilizando herramientas de análisis de redes génicas, con el propósito de comprender de manera preliminar qué rutas moleculares podrían estar implicadas en la regulación de la homeostasis hídrica y el control de la sed.

\subsection{Objetivos específicos}

\begin{enumerate}
    \item Buscar y reunir los genes asociados al fenotipo \textit{Abnormal drinking behaviour} en la base de datos \textbf{Human Phenotype Ontology} (HPO).
    \item Construir una red de interacción entre esos genes empleando herramientas bioinformáticas como \textit{STRING}.
    \item Analizar la red para identificar los genes más conectados (\textit{hubs}) y posibles grupos funcionales.
    \item Explorar los procesos biológicos y rutas moleculares asociados a estos genes mediante análisis de enriquecimiento funcional.
    \item Interpretar los resultados y contrastarlos con la literatura científica, con el fin de proponer hipótesis sobre cómo estos genes podrían influir en los comportamientos anómalos de ingesta de líquidos.
\end{enumerate}

